\scanpage{189}
\mspage{II.37}
% archive sheet number: 188
% block: upper third, lines 1-9 | text
peut surprendre, car (\unsure{en admettant|mettant} que les $\widetilde G_E$ soient semi-simples disons) \deleted{\unsure{on ne voit pas qu'ils soient}} ils ne sont pas nécessairement simplement connexes (au sens de Chevalley) i.e. ils pourraient avoir (a priori\ldots) moins de représentations que leurs Algèbres de Lie. L'explication provient du fait que dans le système projectif des $\widetilde G_{E'}$ \unsure{majorant|dominant} $\widetilde G_E$, \unsure{on a|on trouve} aussi le ``revêtement simplement connexe'' de $\widetilde G_E$, au sens de Chevalley.

% block: middle third, lines 10-16 | text
\deleted{Pour}\inserted{\unsure{on}} s'en \unsure{assure|assurer} directement, \deleted{notons que,} \deleted{\unsure{soit} $\pi_E=\unsure{e|a}$}\inserted{soit} $G'$ le rev. simplement connexe de $\widetilde G_E^0$; et soit $\Gamma$ le noyau de $\deleted{G'\longrightarrow\widetilde G_E}G'(\mathbb Q_\ell)\unsure{\text{---}|\longrightarrow}\widetilde G_E^0(\mathbb Q_\ell)$. C'est un sous-groupe discret de $G'$, donc il existe un sous-groupe ouvert $U$ de $G'(\mathbb Q_\ell)$ tel que $U\longrightarrow\widetilde G_E^0(\mathbb Q_\ell)$ soit une \deleted{\unsure{isom}} application injective (\unsure{i.e. ouverte}). Si $E'$ est le module d'une représentation \inserted{alg} linéaire fidèle de $G'$, il induit une représentation \inserted{analytique} de $U$, donc une représentation d'un $\widetilde G_{E'}^0$\ldots

% block: middle left margin | margin note
% slanted
\margin{laisser tomber}

% block: lower third, lines 17-23 | text
6.8. Si $G$ est un groupe algébrique \deleted{affine}\inserted{\unsure{algébrique|quelconque|linéaire} (\unsure{opérant ds|plutôt de} $\unsure{E|G}$)} défini sur $\mathbb Q_\ell$, alors on peut considérer le groupe analytique $\deleted{\unsure{\pi}}=G(\mathbb Q_\ell)$; \deleted{et le groupe analytique profini \unsure{associé}}\inserted{\unsure{soit} $\pi$ un s-groupe ouvert, \uline{compact}}. Alors \deleted{\unsure{on peut considérer} $\pi$ \unsure{qui est}} le groupe pro-algébrique \inserted{$\pi^{\mathrm{alg}}$} associé est une \uline{extension} de $\pi=\deleted{\unsure{\widetilde\pi}}\unsure{\pi|\overline\pi}$ (considéré comme groupe pro-fini\inserted{-alg}) par
